\documentclass[10pt]{article}

\usepackage[margin=0.5in,landscape]{geometry}
\usepackage{multicol}
\usepackage{amsmath,amssymb}
\usepackage{enumitem}
\setlist{noitemsep}

\begin{document}

\begin{center}
{\Large \textbf{CSCE 222 Exam 1 Cheatsheet}}
\end{center}

\begin{multicols}{2}

%%%%%%%%%%%%%%%%%%%%%%%%%%%%%%%%%%%%%%
\section*{1. Propositional Logic}

\textbf{Proposition}: A statement that is either TRUE or FALSE.

\textbf{Logical Operators}

\begin{tabular}{c|c}
Symbol & Meaning \\
\hline
$\neg p$ & NOT p \\
$p \land q$ & AND \\
$p \lor q$ & OR \\
$p \rightarrow q$ & Implication \\
$p \leftrightarrow q$ & Biconditional
\end{tabular}

\subsection*{Implication Equivalence}

\[
p \rightarrow q \equiv \neg p \lor q
\]

\subsection*{Truth Table Essentials}

\begin{tabular}{c c | c}
p & q & $p \rightarrow q$ \\
\hline
T & T & T \\
T & F & F \\
F & T & T \\
F & F & T
\end{tabular}

\subsection*{Logical Equivalences}

\textbf{Identity}
\[
p \land T = p
\]

\[
p \lor F = p
\]

\textbf{Domination}

\[
p \lor T = T
\]

\[
p \land F = F
\]

\textbf{Idempotent}

\[
p \lor p = p
\]

\[
p \land p = p
\]

\textbf{Double Negation}

\[
\neg(\neg p) = p
\]

\textbf{DeMorgan's Laws}

\[
\neg(p \land q) = \neg p \lor \neg q
\]

\[
\neg(p \lor q) = \neg p \land \neg q
\]

\textbf{Commutative}

\[
p \land q = q \land p
\]

\[
p \lor q = q \lor p
\]

\textbf{Associative}

\[
(p \land q) \land r = p \land (q \land r)
\]

\textbf{Distributive}

\[
p \land (q \lor r) = (p \land q) \lor (p \land r)
\]

\[
p \lor (q \land r) = (p \lor q) \land (p \lor r)
\]

%%%%%%%%%%%%%%%%%%%%%%%%%%%%%%%%%%%%%%
\section*{2. Predicate Logic}

\textbf{Predicate}: A statement with variables. :contentReference[oaicite:0]{index=0}

Example:

\[
P(x): x > 3
\]

\textbf{Quantifiers}

\begin{tabular}{c|c}
Symbol & Meaning \\
\hline
$\forall x$ & For all x (Universal) \\
$\exists x$ & There exists x (Existential)
\end{tabular}

\[
\forall x P(x)
\]

True if P(x) is true for every element in the domain. :contentReference[oaicite:1]{index=1}

\[
\exists x P(x)
\]

True if P(x) is true for at least one element.

\subsection*{Negation of Quantifiers}

\[
\neg(\forall x P(x)) \equiv \exists x \neg P(x)
\]

\[
\neg(\exists x P(x)) \equiv \forall x \neg P(x)
\]

\subsection*{Distribution Rules}

\[
\forall x (P(x) \land Q(x)) \equiv \forall x P(x) \land \forall x Q(x)
\]

\[
\exists x (P(x) \lor Q(x)) \equiv \exists x P(x) \lor \exists x Q(x)
\]

\textbf{Not valid:}

\[
\forall x (P(x) \lor Q(x))
\]

\[
\exists x (P(x) \land Q(x))
\]

\subsection*{Nested Quantifiers}

Order matters:

\[
\forall x \exists y P(x,y) \neq \exists x \forall y P(x,y)
\]

%%%%%%%%%%%%%%%%%%%%%%%%%%%%%%%%%%%%%%
\section*{3. Proof Techniques}

\textbf{Direct Proof}

Assume $p$ and prove $q$.

\[
p \rightarrow q
\]

\textbf{Contrapositive}

\[
p \rightarrow q \equiv \neg q \rightarrow \neg p
\]

\textbf{Proof by Contradiction}

Assume the opposite and derive contradiction.

\[
p \land \neg p
\]

\textbf{Proof by Cases}

Split problem into exhaustive cases.

\textbf{Counterexample}

To disprove:

\[
\forall x P(x)
\]

Find one x where P(x) is false.

%%%%%%%%%%%%%%%%%%%%%%%%%%%%%%%%%%%%%%
\section*{4. Sets}

\textbf{Set Notation}

Roster:

\[
A = \{1,2,3\}
\]

Set-builder:

\[
A = \{x \mid x > 0\}
\]

\textbf{Subset}

\[
A \subseteq B
\]

\[
\forall x (x \in A \rightarrow x \in B)
\]

\textbf{Proper Subset}

\[
A \subset B
\]

A subset but not equal.

\subsection*{Power Set}

\[
P(S) = \text{set of all subsets}
\]

If

\[
|S| = n
\]

then

\[
|P(S)| = 2^n
\]

:contentReference[oaicite:2]{index=2}

\subsection*{Set Operations}

\begin{tabular}{c|c}
Operation & Meaning \\
\hline
$A \cup B$ & Union \\
$A \cap B$ & Intersection \\
$A-B$ & Difference \\
$A^c$ & Complement
\end{tabular}

\subsection*{Formal Definitions}

\[
A \cup B = \{x \mid x \in A \lor x \in B\}
\]

\[
A \cap B = \{x \mid x \in A \land x \in B\}
\]

\[
B-A = \{x \mid x \in B \land x \notin A\}
\]

\subsection*{Cartesian Product}

\[
A \times B = \{(a,b) \mid a \in A, b \in B\}
\]

If

\[
|A| = n, |B| = m
\]

\[
|A \times B| = nm
\]

\subsection*{Relation}

\[
R \subseteq A \times B
\]

%%%%%%%%%%%%%%%%%%%%%%%%%%%%%%%%%%%%%%
\section*{5. Functions}

\textbf{Function}

\[
f : A \rightarrow B
\]

Maps every element of A to exactly one element in B. :contentReference[oaicite:3]{index=3}

\subsection*{Terminology}

\begin{tabular}{c|c}
Term & Meaning \\
\hline
Domain & Input set \\
Codomain & Output set \\
Range & Actual outputs \\
Image & $f(x)$ \\
Preimage & x such that $f(x)=y$
\end{tabular}

\subsection*{One-to-One (Injective)}

\[
f(x_1)=f(x_2) \Rightarrow x_1=x_2
\]

Each output has at most one input. :contentReference[oaicite:4]{index=4}

\subsection*{Onto (Surjective)}

\[
\forall y \in Y, \exists x \in X: f(x)=y
\]

Every output is used.

\subsection*{Bijection}

Function is both:

\begin{itemize}
\item Injective
\item Surjective
\end{itemize}

\subsection*{Inverse Function}

Exists only if function is bijective.

\[
f^{-1}(f(x)) = x
\]

\subsection*{Floor and Ceiling}

\[
\lfloor x \rfloor
\]

Largest integer $\le x$.

\[
\lceil x \rceil
\]

Smallest integer $\ge x$.

%%%%%%%%%%%%%%%%%%%%%%%%%%%%%%%%%%%%%%
\section*{6. Algorithms}

\textbf{Algorithm}: Finite sequence of instructions solving a problem.

\subsection*{Complexity Types}

\begin{tabular}{c|c}
Notation & Meaning \\
\hline
$O(1)$ & Constant \\
$O(\log n)$ & Logarithmic \\
$O(n)$ & Linear \\
$O(n^2)$ & Quadratic \\
$O(n^b)$ & Polynomial \\
$O(b^n)$ & Exponential
\end{tabular}

:contentReference[oaicite:5]{index=5}

\subsection*{Time Complexity}

Number of operations executed.

Simplifying assumption: all operations cost equal time. :contentReference[oaicite:6]{index=6}

\subsection*{Common Algorithm Complexities}

\begin{tabular}{c|c|c}
Algorithm & Best & Worst \\
\hline
Linear Search & $\Omega(1)$ & $O(n)$ \\
Binary Search & $\Omega(1)$ & $O(\log n)$ \\
Selection Sort & $\Omega(n^2)$ & $O(n^2)$ \\
Insertion Sort & $\Omega(n)$ & $O(n^2)$ \\
Merge Sort & $\Omega(n\log n)$ & $O(n\log n)$
\end{tabular}

%%%%%%%%%%%%%%%%%%%%%%%%%%%%%%%%%%%%%%
\section*{7. Asymptotic Notation}

\subsection*{Big-O}

Upper bound.

\[
f(n) \in O(g(n))
\]

if

\[
f(n) \le Cg(n)
\]

for large n.

\subsection*{Big-Omega}

Lower bound.

\[
f(n) \in \Omega(g(n))
\]

if

\[
f(n) \ge Cg(n)
\]

:contentReference[oaicite:7]{index=7}

\subsection*{Big-Theta}

Tight bound.

\[
f(n) \in \Theta(g(n))
\]

if

\[
f(n) \in O(g(n)) \text{ and } f(n) \in \Omega(g(n))
\]

%%%%%%%%%%%%%%%%%%%%%%%%%%%%%%%%%%%%%%
%%%%%%%%%%%%%%%%%%%%%%%%%%%%%%%%%%%%%%
\section*{8. Common Proof Approaches (Exam Patterns)}

\subsection*{How to Choose a Proof Technique}

\begin{tabular}{l|l}
Problem Form & Technique \\
\hline
$\forall x \; P(x) \rightarrow Q(x)$ & Direct proof or Contrapositive \\
$\neg(\forall x P(x))$ & Find counterexample \\
$\exists x P(x)$ & Construct example \\
Set equality $A=B$ & Show $A\subseteq B$ and $B\subseteq A$ \\
Function injective & Assume $f(x_1)=f(x_2)$ \\
Function surjective & Solve $y=f(x)$ for $x$ \\
Implication false & Counterexample
\end{tabular}

%%%%%%%%%%%%%%%%%%%%%%%%%%%%%%%%%%%%%%
\subsection*{Direct Proof Template}

Used when proving implications.

Goal:

\[
P \rightarrow Q
\]

Steps:

\begin{enumerate}
\item Assume $P$ is true
\item Use definitions / algebra / known facts
\item Derive $Q$
\item Conclude $P \rightarrow Q$
\end{enumerate}

Example structure:

\[
\text{Assume } x \in A
\]

\[
\text{Show } x \in B
\]

%%%%%%%%%%%%%%%%%%%%%%%%%%%%%%%%%%%%%%
\subsection*{Contrapositive Proof}

Useful when $Q$ is easier to negate.

\[
P \rightarrow Q
\]

Equivalent statement:

\[
\neg Q \rightarrow \neg P
\]

Proof steps:

\begin{enumerate}
\item Assume $\neg Q$
\item Show $\neg P$
\end{enumerate}

Example:

\[
\text{If } n^2 \text{ is even } \rightarrow n \text{ is even}
\]

Prove instead:

\[
n \text{ is odd } \rightarrow n^2 \text{ is odd}
\]

%%%%%%%%%%%%%%%%%%%%%%%%%%%%%%%%%%%%%%
\subsection*{Proof by Contradiction}

Used when direct proof is difficult.

Steps:

\begin{enumerate}
\item Assume the statement is false
\item Derive contradiction
\item Conclude statement must be true
\end{enumerate}

Contradiction form:

\[
P \land \neg P
\]

Example:

\[
\sqrt{2} \text{ is irrational}
\]

%%%%%%%%%%%%%%%%%%%%%%%%%%%%%%%%%%%%%%
\subsection*{Proof by Cases}

Used when domain splits naturally.

Example:

\[
n \text{ is even OR odd}
\]

Structure:

Case 1: $n$ even

\[
n = 2k
\]

Case 2: $n$ odd

\[
n = 2k+1
\]

Show both cases satisfy the result.

%%%%%%%%%%%%%%%%%%%%%%%%%%%%%%%%%%%%%%
\subsection*{Counterexample}

Used to disprove universal statements.

Statement:

\[
\forall x P(x)
\]

Find one element:

\[
x=a
\]

such that:

\[
P(a) \text{ is false}
\]

Example:

\[
\forall x (x^2 > x)
\]

Counterexample:

\[
x=0
\]

%%%%%%%%%%%%%%%%%%%%%%%%%%%%%%%%%%%%%%
\section*{9. Common Set Proof Patterns}

\subsection*{Subset Proof}

To prove:

\[
A \subseteq B
\]

Steps:

\begin{enumerate}
\item Let $x \in A$
\item Use definitions
\item Show $x \in B$
\end{enumerate}

Example structure:

\[
x \in A
\]

\[
\Rightarrow x \text{ satisfies property of B}
\]

\[
\Rightarrow x \in B
\]

%%%%%%%%%%%%%%%%%%%%%%%%%%%%%%%%%%%%%%
\subsection*{Set Equality}

To prove:

\[
A = B
\]

Must prove two subsets:

\[
A \subseteq B
\]

\[
B \subseteq A
\]

Structure:

\[
x \in A \Rightarrow x \in B
\]

\[
x \in B \Rightarrow x \in A
\]

%%%%%%%%%%%%%%%%%%%%%%%%%%%%%%%%%%%%%%
\subsection*{Disjoint Sets}

Goal:

\[
A \cap B = \emptyset
\]

Method:

Assume

\[
x \in A \cap B
\]

Show contradiction.

%%%%%%%%%%%%%%%%%%%%%%%%%%%%%%%%%%%%%%
\section*{10. Function Proof Patterns}

\subsection*{Proving Injective (One-to-One)}

Goal:

\[
f(x_1)=f(x_2) \Rightarrow x_1=x_2
\]

Steps:

\begin{enumerate}
\item Assume $f(x_1)=f(x_2)$
\item Substitute definition of $f$
\item Solve algebraically
\item Show $x_1=x_2$
\end{enumerate}

Example:

\[
f(x)=3x+5
\]

\[
3x_1+5 = 3x_2+5
\]

\[
3x_1 = 3x_2
\]

\[
x_1=x_2
\]

%%%%%%%%%%%%%%%%%%%%%%%%%%%%%%%%%%%%%%
\subsection*{Proving Surjective (Onto)}

Goal:

\[
\forall y \in B, \exists x \in A
\]

such that

\[
f(x)=y
\]

Steps:

\begin{enumerate}
\item Let $y \in B$
\item Solve $y=f(x)$ for $x$
\item Show $x$ is in domain
\end{enumerate}

%%%%%%%%%%%%%%%%%%%%%%%%%%%%%%%%%%%%%%
\subsection*{Proving Function Equality}

To prove:

\[
f = g
\]

Must show:

\[
f(x)=g(x) \quad \forall x
\]

%%%%%%%%%%%%%%%%%%%%%%%%%%%%%%%%%%%%%%
\section*{11. Logic Argument Proof Pattern}

For validity proofs:

\begin{enumerate}
\item Write premises
\item Apply rules of inference
\item Derive conclusion
\end{enumerate}

Common inference rules:

\begin{tabular}{l}
Modus Ponens \\
Modus Tollens \\
Hypothetical Syllogism \\
Disjunctive Syllogism \\
Resolution
\end{tabular}

%%%%%%%%%%%%%%%%%%%%%%%%%%%%%%%%%%%%%%
\end{multicols}

\end{document}